\documentclass[../../main.tex]{subfiles}
\graphicspath{{../../image/}{../../figure/}}
\begin{document}

\section{等熵气体动力学方程组}\label{section:ch17-等熵气体动力学}

理想气体等熵气体动力学方程组\eqref{align:ch07-理想气体}在 Euler 坐标下，是有史以来第一个被导出的双曲守恒律方程组，
过去两个世纪一直充当检验该理论的试验场。因此，用补偿紧致性方法应用于该方程组来结束本章是适宜的。

有启发意义的是同时在其原始形式\eqref{align:ch07-理想气体}（状态变量为密度 $\rho$ 与速度 $v$）
与其典则形式
\begin{align}
\begin{cases}
\partial_t\rho+\partial_x m=0\\
\partial_t m+\partial_x\Bigl[\dfrac{m^2}{\rho}+\kappa\rho^\gamma\Bigr]=0,
\end{cases}
\label{align:ch17-1781}
\end{align}
（状态变量为密度 $\rho$ 与动量 $m=\rho v$）中监测该方程组。密度的物理范围是 $0\leqslant\rho<\infty$，而 $v$ 与 $m$ 可取 $(-\infty,\infty)$ 中任意值。

为方便起见，我们把状态变量标度化，使 $\kappa=(\gamma-1)^2/(4\gamma)$，并采用记号 $\theta=\frac{1}{2}(\gamma-1)$，此时特征速度\eqref{align:ch07-理想气体特征速度}与 Riemann 不变量\eqref{align:ch07-理想气体Riemann}取形式
\begin{align}
\lambda=-\theta\rho^\theta+v=-\theta\rho^\theta+\frac{m}{\rho},\qquad\mu=\theta\rho^\theta+v=\theta\rho^\theta+\frac{m}{\rho},
\label{align:ch17-1782}
\end{align}
\begin{align}
z=-\rho^\theta+v=-\rho^\theta+\frac{m}{\rho},\qquad w=\rho^\theta+v=\rho^\theta+\frac{m}{\rho}.
\label{align:ch17-1783}
\end{align}

不难构造取值于状态空间 $[0,\infty)\times(-\infty,\infty)$ 紧集的近似解序列。例如，可沿相对于方程组
\begin{align}
\begin{cases}
\partial_t\rho+\partial_x m=\mu\,\partial_x^2\rho\\
\partial_t m+\partial_x\Bigl[\dfrac{m^2}{\rho}+\kappa\rho^\gamma\Bigr]=\mu\,\partial_x^2m
\end{cases}
\label{align:ch17-1784}
\end{align}
的消失黏性方法前进，该方程组允许（正）不变区域族
\begin{align}
\mathscr{U}_M=\{(\rho,m):\rho\geqslant0,\ -M\leqslant z(\rho,m)\leqslant w(\rho,m)\leqslant M\}.
\label{align:ch17-1785}
\end{align}
此外，\eqref{align:ch17-1784}在上半平面上以含于 $L^\infty(-\infty,\infty)\cap L^2(-\infty,\infty)$ 有界集的初值为初值的解 $(\rho_\mu,m_\mu)$，满足
\begin{align}
\partial_t\eta(\rho_\mu,m_\mu)+\partial_x q(\rho_\mu,m_\mu)\subset\text{紧集于 }W^{-1,2}_{\mathrm{loc}},
\label{align:ch17-1786}
\end{align}
对\eqref{align:ch17-1781}的任意熵--熵流对 $(\eta,q)$ 成立。具有类似性质的近似解也由有限差分格式构造，如 Lax--Friedrichs 格式与 Godunov 格式。它们都导出下述存在性定理：

\begin{theorem}[等熵气体动力学的存在性]\label{theorem:ch17-气体存在性}
对任意 $\gamma>1$，存在方程组\eqref{align:ch07-理想气体}在上半平面上的有界解 $(\rho,v)$，
其指定初值为
\begin{align}
(\rho(x,0),v(x,0))=(\rho_0(x),v_0(x)),\qquad-\infty<x<\infty,
\label{align:ch17-1787}
\end{align}
其中 $(\rho_0,v_0)\in L^\infty(-\infty,\infty)$ 且 $\rho_0(x)\geqslant0$（$-\infty<x<\infty$）。此外，该解满足熵可容性条件
\begin{align}
\partial_t\eta(\rho,m)+\partial_x q(\rho,m)\leqslant0,
\label{align:ch17-1788}
\end{align}
对\eqref{align:ch17-1781}的任意熵--熵流对 $(\eta,q)$（其中 $\eta(\rho,m)$ 凸）成立。
\end{theorem}

证明使用\eqref{align:ch17-17310}来确立：与近似解序列相联系的 Young 测度的支撑，要么约化为状态空间中一点，要么被限制在轴 $\rho=0$（真空态）上。

作为 $(\rho,v)$ 的函数，\eqref{align:ch07-理想气体}的任意熵 $\eta$ 满足可积性条件
\begin{align}
\eta_{\rho\rho}=\theta^2\rho^{\gamma-3}\eta_{vv}.
\label{align:ch17-1789}
\end{align}
上述方程沿轴 $\rho=0$ 奇异，且当越过阈值 $\gamma=3$ 时奇性的本质发生变化。因此，处理 $\gamma<3$ 与 $\gamma>3$ 的情形须用不同论证。

与这里相关的是所谓弱熵--熵流对，它们在 $\rho=0$ 处为零。它们容许表示
\begin{align}
\begin{cases}
\eta(\rho,v)=\displaystyle\int_{-\infty}^\infty\chi(\rho,\omega-v)g(\omega)\,d\omega\\
q(\rho,v)=\displaystyle\int_{-\infty}^\infty\chi(\rho,\omega-v)[\theta\omega+(1-\theta)v]g(\omega)\,d\omega,
\end{cases}
\label{align:ch17-17810}
\end{align}
其中
\begin{align}
\chi(\rho,v)=
\begin{cases}
(\rho^{2\theta}-v^2)^s,&\text{若 }\rho^{2\theta}>v^2,\\
0,&\text{若 }\rho^{2\theta}\leqslant v^2,
\end{cases}
\label{align:ch17-17811}
\end{align}
其中 $s=\frac{1}{2}\frac{3-\gamma}{\gamma-1}$。于是 $\chi$ 是\eqref{align:ch17-1789}在初值 $\eta(0,v)=0$、$\eta_\rho(0,v)=\delta_0(v)$ 下的基本解。

如\cref{section:ch02-热弹性}已指出，经典分子动理论预言：对具有 $n$ 个自由度的气体，
绝热指数取值 $\gamma=1+\frac{2}{n}$。当自由度数为奇数 $n=2\ell+1$ 时，\eqref{align:ch17-17811}中的指数 $s$ 为整数 $\ell$。在这一特殊情形下，弱熵的分析——从而 Young 测度的约化——大大简化。然而，即使在这种较简单的情形下，证明也相当技术化，故将其归入\cref{section:ch17-notes}所引文献。这里只展示退化情形 $\gamma=3$。

当 $\gamma=3$（即 $\theta=1$）时，\eqref{align:ch17-1782}与\eqref{align:ch17-1783}给出 $\lambda=z$、$\mu=w$，此时两特征族完全解耦。特别地，\eqref{align:ch12-ent流}退化为 $q_z=z\eta_z$、$q_w=w\eta_w$，
从而存在只依赖于 $z$ 的熵--熵流对 $(\eta,q)$，例如 $(2z,z^2)$ 与 $(3z^2,2z^3)$。

现设\eqref{align:ch17-1784}的解序列 $\{(\rho_{\mu_k},m_{\mu_k})\}$（$\mu_k\to0$，$k\to\infty$）诱导出弱收敛的 Riemann 不变量子列 $\{(z_j,w_j)\}$，相联系的 Young 测度族为 $\nu_{x,t}$。固定 $(x,t)$，令 $\nu_{x,t}=\nu$，并对两个熵--熵流对 $(2z,z^2)$ 与 $(3z^2,2z^3)$ 应用\eqref{align:ch17-17310}，得
\begin{align}
4\langle\nu,z\rangle\langle\nu,z^3\rangle-3\langle\nu,z^2\rangle\langle\nu,z^2\rangle=\langle\nu,z^4\rangle.
\label{align:ch17-17812}
\end{align}

接着考虑不等式
\begin{align}
z^4-4z^3\bar z+6z^2\bar z^2-4z\bar z^3+\bar z^4=(z-\bar z)^4\geqslant0,
\label{align:ch17-17813}
\end{align}
其中 $\bar z=\langle\nu,z\rangle$，并把测度 $\nu$ 作用于它，得
\begin{align}
\langle\nu,z^4\rangle-4\langle\nu,z^3\rangle\langle\nu,z\rangle+6\langle\nu,z^2\rangle\langle\nu,z\rangle^2-3\langle\nu,z\rangle^4\geqslant0.
\label{align:ch17-17814}
\end{align}
将\eqref{align:ch17-17814}与\eqref{align:ch17-17812}结合，得
\begin{align}
-3\bigl[\langle\nu,z^2\rangle-\langle\nu,z\rangle^2\bigr]^2\geqslant0,
\label{align:ch17-17815}
\end{align}
由此 $\langle\nu,z^2\rangle=\langle\nu,z\rangle^2$。因此 $\{z_j\}$ 强收敛于 $\bar z=\langle\nu,z\rangle$。

类似地可证 $\{w_j\}$ 强收敛于 $\bar w=\langle\nu,w\rangle$。特别地，$(\bar z,\bar w)$ 通过 $\bar\rho=\frac{1}{2}(\bar w-\bar z)$、$\bar v=\frac{1}{2}(\bar w+\bar z)$ 诱导出\eqref{align:ch07-理想气体}的解 $(\bar\rho,\bar v)$。

从连续介质物理的立场看，更自然的是采用物理的而非人工的消失黏性，通过方程组
\begin{align}
\begin{cases}
\partial_t\rho+\partial_x m=0\\
\partial_t m+\partial_x\Bigl[\dfrac{m^2}{\rho}+\kappa\rho^\gamma\Bigr]=\mu\,\partial_x^2v
\end{cases}
\label{align:ch17-17816}
\end{align}
（代替\eqref{align:ch17-1784}）构造\eqref{align:ch17-1781}的近似解 $(\rho_\mu,m_\mu)$。然而，这因需要\eqref{align:ch17-17816}解的 $L^\infty$ 界（对 $\mu>0$ 一致成立）而受阻，而这类界目前尚不可得。人们可尝试改用 $L^p$ 界，但此时把 Young 测度约化为点质量的任务变得相当费力。下面给出沿此方向所建立的更一般结果的一个简化版本。它涉及当 $|x|\to\infty$ 时趋于状态 $(\bar\rho,0)$（$\bar\rho>0$）的初值 $(\rho_0,m_0)$。

\begin{theorem}[等熵气体的物理黏性极限]\label{theorem:ch17-气体物理黏性}
考虑光滑初值族 $\{(\rho_{0\mu},m_{0\mu}),\ \mu>0\}$，具有下述性质：密度为正，
\begin{align}
\rho_{0\mu}(x)\geqslant a_\mu>0,\qquad-\infty<x<\infty,
\label{align:ch17-17817}
\end{align}
且其振荡受控于
\begin{align}
\mu^2\int_{-\infty}^\infty\frac{[\partial_x\rho_{0\mu}(x)]^2}{\rho_{0\mu}^3(x)}\,dx\leqslant N.
\label{align:ch17-17818}
\end{align}
动量一致有界，
\begin{align}
\int_{-\infty}^\infty|m_{0\mu}(x)|\,dx\leqslant P,
\label{align:ch17-17819}
\end{align}
且相对于状态 $(\bar\rho,0)$ 的机械能也一致有界：
\begin{align}
\int_{-\infty}^\infty\Bigl(\frac{\kappa}{\gamma-1}[\rho_{0\mu}^\gamma(x)-\bar\rho^\gamma-\gamma\bar\rho^{\gamma-1}(\rho_{0\mu}(x)-\bar\rho)]+\frac{1}{2}\frac{m_{0\mu}^2(x)}{\rho_{0\mu}(x)}\Bigr)\,dx\leqslant E.
\label{align:ch17-17820}
\end{align}
最后，
\begin{align}
(\rho_{0\mu},m_{0\mu})\longrightarrow(\rho_0,m_0),\qquad\mu\to0
\label{align:ch17-17821}
\end{align}
在分布意义下成立。则对任意 $\mu>0$，存在以 $(\rho_{0\mu},m_{0\mu})$ 为初值的\eqref{align:ch17-17816}Cauchy 问题的古典解 $(\rho_\mu,m_\mu)$。此外，存在序列 $\{\mu_j\}$（$\mu_j\to0$，$j\to\infty$），使得 $\{(\rho_{\mu_j},m_{\mu_j})\}$ 几乎处处收敛于以 $(\rho_0,m_0)$ 为初值的\eqref{align:ch17-1781}Cauchy 问题的解 $(\rho,m)$。该解满足熵可容性条件。
\end{theorem}

上述命题的冗长而技术化的证明（见\cref{section:ch17-notes}所引文献）首先确立\eqref{align:ch17-17816}解 $(\rho_\mu,m_\mu)$ 的 $L^p$ 界（关于 $\mu>0$ 一致），然后证明诱导的 Young 测度约化为点质量，从而 $(\rho_\mu,m_\mu)$ 的 $\mu\to0$ 弱极限 $(\rho,m)$ 是\eqref{align:ch17-1781}的解。为传达证明方法论的味道，下面勾勒先验估计的推导。不过，为简化陈述，我们将在外加假设 $\bar\rho=0$、且初值总质量有限且一致有界下确立估计：
\begin{align}
\int_{-\infty}^\infty\rho_{0\mu}\,dx\leqslant M.
\label{align:ch17-17822}
\end{align}
尽管这些要求违反定理的条件，但就估计的推导而言，它们并不影响论证的本质。

由质量与机械能的平衡律，结合\eqref{align:ch17-17822}与\eqref{align:ch17-17820}，得
\begin{align}
\int_{-\infty}^\infty\rho_\mu(x,t)\,dx\leqslant M,
\label{align:ch17-17823}
\end{align}
\begin{align}
\int_{-\infty}^\infty\Bigl(\frac{\kappa}{\gamma-1}\rho_\mu^\gamma(x,t)+\frac{1}{2}\frac{m_\mu^2(x,t)}{\rho_\mu(x,t)}\Bigr)\,dx+\mu\int_0^t\int_{-\infty}^\infty[\partial_xv_\mu(x,\tau)]^2\,dx\,d\tau\leqslant E.
\label{align:ch17-17824}
\end{align}

接着引入函数
\begin{align}
\Phi_\mu(x,t)=\int_{-\infty}^x m_\mu(y,t)\,dy,
\label{align:ch17-17825}
\end{align}
它充当一个势函数；注意到由\eqref{align:ch17-17823}、\eqref{align:ch17-17824}与 Schwarz 不等式，我们有
\begin{align}
|\Phi_\mu(x,t)|^2\leqslant\int_{-\infty}^\infty\frac{m_\mu^2(x,t)}{\rho_\mu(x,t)}\,dx\int_{-\infty}^\infty\rho_\mu(x,t)\,dx\leqslant 2EM.
\label{align:ch17-17826}
\end{align}
由\eqref{align:ch17-17816}与\eqref{align:ch17-17825}，
\begin{align}
\partial_t(\rho_\mu\Phi_\mu)+\partial_x(m_\mu\Phi_\mu)=-\kappa\rho_\mu^{\gamma+1}+\mu\rho_\mu\partial_xv_\mu.
\label{align:ch17-17827}
\end{align}
将\eqref{align:ch17-17827}在 $(-\infty,\infty)\times(0,t)$ 上积分，使用\eqref{align:ch17-17823}、\eqref{align:ch17-17824}、\eqref{align:ch17-17826}，并设 $\mu\leqslant\kappa$，我们断言
\begin{align}
\int_0^t\int_{-\infty}^\infty\rho_\mu^{\gamma+1}(x,\tau)\,dx\,d\tau\leqslant\frac{2}{\kappa}E+\frac{8}{\kappa}M^3+\frac{\mu t}{\kappa}M,
\label{align:ch17-17828}
\end{align}
这就提供了第一个基本新估计。

为得到下一个估计，我们切换到 Lagrange 坐标 $(\xi,t)$。为避免混淆，Lagrange 坐标中关于 $\xi$ 与 $t$ 的偏导数以下标 $\xi$、$t$ 表示，而符号 $\partial_x$、$\partial_t$ 保留给 Euler 坐标中关于 $x$ 与 $t$ 的偏导数。特别地，$x_\xi=\rho^{-1}$、$x_t=v$、$\partial_x\xi=\rho$。

在 Lagrange 坐标下，方程组\eqref{align:ch17-17816}读作
\begin{align}
\begin{cases}
(\rho^{-1})_t-v_\xi=0\\
v_t+(\kappa\rho^\gamma)_\xi=\mu(\rho v_\xi)_\xi.
\end{cases}
\label{align:ch17-17829}
\end{align}

直接计算并使用\eqref{align:ch17-17829}$_1$，得
\begin{align}
[(\partial_x\rho_\mu^{-1})^2]_t=[\rho_\mu^{-2}(\rho_{\mu\xi}v_\mu)^2]_t=-2\rho_\mu^{-1}\rho_{\mu\xi}(\rho_\mu v_{\mu\xi})_\xi.
\label{align:ch17-17830}
\end{align}
将\eqref{align:ch17-17830}乘以 $\mu$，再把项 $\mu(\rho_\mu v_{\mu\xi})_\xi$ 由\eqref{align:ch17-17829}$_2$ 代入所得方程，我们得到
\begin{align}
\begin{split}
\mu[(\partial_x\rho_\mu^{-1})^2]_t&=-2\rho_\mu^{-1}\rho_{\mu\xi}v_{\mu t}-2\kappa\gamma\rho_\mu^{\gamma-2}(\rho_{\mu\xi})^2\\
&=-2(\rho_\mu^{-1}\rho_{\mu\xi}v_\mu)_t-2(\rho_\mu v_\mu v_{\mu\xi})_\xi+2\rho_\mu(v_\mu v_{\mu\xi})^2-2\kappa\gamma\rho_\mu^{\gamma-2}(\rho_{\mu\xi})^2.
\end{split}
\label{align:ch17-17831}
\end{align}

现对\eqref{align:ch17-17831}先关于 $\xi$ 在 $(-\infty,\infty)$ 上积分，再关于时间在 $(0,t)$ 上积分。最后切换回 Euler 坐标 $(x,t)$。最终结果是下述恒等式：
\begin{align}
\begin{split}
\mu\int_{-\infty}^\infty&\rho_\mu^{-3}(\partial_x\rho_\mu)^2\,dx+2\kappa\gamma\int_0^t\int_{-\infty}^\infty\rho_\mu^{\gamma-3}(\partial_x\rho_\mu)^2\,dx\,d\tau\\
=-2&\int_{-\infty}^\infty\rho_\mu^{-1}(\partial_x\rho_\mu)v_\mu\,dx+2\int_0^t\int_{-\infty}^\infty(\partial_xv_\mu)^2\,dx\,d\tau\\
+\mu&\int_{-\infty}^\infty\rho_{0\mu}^{-3}(\partial_x\rho_{0\mu})^2\,dx+2\int_{-\infty}^\infty\rho_{0\mu}^{-1}(\partial_x\rho_{0\mu})v_{0\mu}\,dx.
\end{split}
\label{align:ch17-17832}
\end{align}

我们将\eqref{align:ch17-17832}乘以 $\mu$，并如下估计右端第一项与最后一项：
\begin{align}
-2\mu\int_{-\infty}^\infty\rho_\mu^{-1}(\partial_x\rho_\mu)v_\mu\,dx\leqslant\frac{\mu^2}{2}\int_{-\infty}^\infty\rho_\mu^{-3}(\partial_x\rho_\mu)^2\,dx+2\int_{-\infty}^\infty\rho_\mu^{-1}m_\mu^2\,dx,
\label{align:ch17-17833}
\end{align}
\begin{align}
2\mu\int_{-\infty}^\infty\rho_{0\mu}^{-1}(\partial_x\rho_{0\mu})v_{0\mu}\,dx\leqslant\frac{\mu^2}{2}\int_{-\infty}^\infty\rho_{0\mu}^{-3}(\partial_x\rho_{0\mu})^2\,dx+2\int_{-\infty}^\infty\rho_{0\mu}^{-1}m_{0m}^2\,dx.
\label{align:ch17-17834}
\end{align}
\eqref{align:ch17-17832}右端第三项由\eqref{align:ch17-17818}估计，% 原书此处误排为 (18.7.32)
而\eqref{align:ch17-17832}、\eqref{align:ch17-17833}、\eqref{align:ch17-17834}右端第二项都借助\eqref{align:ch17-17824}估计。于是我们得到期望的估计：
\begin{align}
\mu^2\int_{-\infty}^\infty\rho_\mu^{-3}(x,t)(\partial_x\rho_\mu(x,t))^2\,dx+4\mu\kappa\gamma\int_0^t\int_{-\infty}^\infty\rho_\mu^{\gamma-3}(\partial_x\rho_\mu)^2\,dx\,d\tau\leqslant20E+3N,
\label{align:ch17-17835}
\end{align}
它控制密度的振荡。

类似但仅是局部的估计甚至当 $\bar\rho>0$ 时也适用。为建立令人满意的紧致性框架，还需要 $\rho$ 与 $v$ 上更高的局部 $L^p$ 界，这可借助弱熵--熵流对\eqref{align:ch17-17810}（取 $g(\omega)=\frac{1}{2}|\omega|\omega$）得到。然而，这些估计的推导关键在于：在上半平面的相当一部分上保持密度远离零，因而要求 $\bar\rho>0$。

凭借上述估计，可证明 Young 测度约化为点质量。然而，该任务技术上富有挑战，因为测度在相空间 $\{(\rho,m):0<\rho<\infty,\ -\infty<m<\infty\}$ 上取值，而该空间非紧。困难由一个富有洞见的相空间紧化所解决。读者可查阅\cref{section:ch17-notes}所引文献。

对一大类含非齐性与/或带奇性源项、并表现出几何效应的气体动力学一维方程组，使用 $L^p$（而非 $L^\infty$）界变得不可或缺。一个有启发性的例子是支配理想气体三维径向等熵流的方程组\eqref{align:ch07-径向流}，
我们在此写成
\begin{align}
\begin{cases}
\partial_t(r^2\rho)+\partial_r(r^2m)=0\\
\partial_t(r^2m)+\partial_r\Bigl[r^2\Bigl(\dfrac{m^2}{\rho}+\kappa\rho^\gamma\Bigr)\Bigr]=2\kappa r\rho^\gamma.
\end{cases}
\label{align:ch17-17836}
\end{align}

只要流动被阻碍而不能趋近原点（例如被一个实心球形障碍物），该方程组就可通过 $L^\infty$ 估计处理。然而，对定义在一切 $r\geqslant0$ 上的流动，密度与速度可能在原点处变为无穷，因此通过 $L^p$ 估计处理是成功的唯一希望。相应地，对 $r\geqslant0$、$t\geqslant0$ 上的 Cauchy 问题，在质量守恒
\begin{align}
\int_0^\infty r^2\rho(r,t)\,dr=M,\qquad 0\leqslant t<\infty
\label{align:ch17-17837}
\end{align}
且能量有限
\begin{align}
\int_0^\infty r^2\Bigl[\frac{\kappa}{\gamma-1}\rho^\gamma(r,t)+\frac{m^2(r,t)}{2\rho(r,t)}\Bigr]\,dr\leqslant E,\qquad 0\leqslant t<\infty
\label{align:ch17-17838}
\end{align}
的函数类中，已建立下述存在性定理：

\begin{theorem}[径向等熵气体的存在性]\label{theorem:ch17-径向气体}
方程组\eqref{align:ch17-17836}（$3/2<\gamma\leqslant3$）以具有有限质量 $M$ 与有限能量 $E$ 的初值为初值的 Cauchy 问题，在 $[0,\infty)\times[0,\infty)$ 上有一个弱解 $(\rho(r,t),m(r,t))$，满足\eqref{align:ch17-17837}与\eqref{align:ch17-17838}。
\end{theorem}

上述命题的证明（见\cref{section:ch17-notes}所引文献）与\cref{theorem:ch17-气体物理黏性}的证明极为相似。特别地，核心作用由估计
\begin{align}
\int_0^t\int_0^\infty r^4\rho^{\gamma+1}(r,\tau)\,dr\,d\tau\leqslant\frac{1}{\kappa}\sqrt{2M^3E},
\label{align:ch17-17839}
\end{align}
即\eqref{align:ch17-17828}的类比，所扮演。为形式地推导\eqref{align:ch17-17839}，注意\eqref{align:ch17-17836}$_1$ 蕴含存在一个 Lipschitz 连续的势函数 $\Psi$：
\begin{align}
\partial_r\Psi=r^2\rho,\qquad\partial_t\Psi=-r^2m.
\label{align:ch17-17840}
\end{align}
结合\eqref{align:ch17-17836}$_2$ 与\eqref{align:ch17-17840}，易验证恒等式
\begin{align}
\partial_t(r^2m\Psi)+\partial_r\Bigl[r^2\Bigl(\frac{m^2}{\rho}+\kappa\rho^\gamma\Bigr)\Psi\Bigr]=\kappa r^4\rho^{\gamma+1}+2\kappa r\rho^\gamma\Psi.
\label{align:ch17-17841}
\end{align}
我们通过 $\Psi(0,t)=0$ 规范化 $\Psi$，此时\eqref{align:ch17-17840}与\eqref{align:ch17-17837}蕴含 $0\leqslant\Psi(r,t)\leqslant M$。于是 Schwarz 不等式连同\eqref{align:ch17-17837}、\eqref{align:ch17-17838}给出
\begin{align}
\Bigl[\int_0^\infty r^2m\Psi\,dr\Bigr]^2\leqslant M^2\int_0^\infty r^2\frac{m^2}{\rho}\,dr\int_0^\infty r^2\rho\,dr\leqslant2M^3E.
\label{align:ch17-17842}
\end{align}
因此，将\eqref{align:ch17-17841}在 $[0,\infty)\times[0,t]$ 上积分，就得到\eqref{align:ch17-17839}。

\end{document}
